\subsection{Counter retrieval: a practical adaptive precision separation}
\label{sec:counter-precision-results}
The counter study tests the same precision objective as the separation
theorem on existing frozen ML comparisons. Its operational decision is
whether to reserve feedback between two register retrievals in order to
meet a required precision for every feasible energy allocation.
It uses the three actual cumulative consumption counters described in
Appendix~\ref{app:historical-design}, their 91 eligible test days, four
existing forecasters and all six pairs. Blocks are $[00,06)$, $[06,12)$
and $[12,18)$ UTC, yielding $3\times91\times3\times6=4,914$ audits.
The 18:00--24:00 block is excluded to avoid ambiguity about the next
day's forecast availability at midnight. Each audit occurs just after its block ends. Its start/end counters and
day-start predictions are permitted. The following day's persistence
forecasts are unavailable under this declared issue schedule; they are
not suppressed from an already completed full-period archive. No measured
arrival times are inferred. Increments use integer subtraction at exact
UTC endpoints before conversion to kWh, with no interpolation.

The interface returns one stored cumulative register per request. Five
internal hourly registers are eligible and the hard budget is two per
standalone pair/block audit, for every policy. The six pairs represent
alternative audits; a joint four-model workflow could share registers and
would require separate cost accounting. Hourly increments and full-trace exports are
not offered at that price. This is a specified access replay, not a claim
that the source deployed such a retrieval service. All information admitted
by this per-audit protocol is used; concentration-profile witnesses satisfy the retained
total and nonnegativity constraints. No empirical maximum is imposed as a
physical cap. Additional defensible caps or measurements would require
recomputing both policies and adverse witnesses.

The protocol was frozen before the new geometry calculations and
audit-answer replay. The primary requirement is $W\leq0.5W_0$: a
scale-free uncertainty-reduction service level known before retrieval,
not a calibrated economic loss threshold or a guarantee of a winning
model. An operational auditor can instead specify an absolute MSE-width
budget; the policy construction accepts that threshold. Here
the fixed sensitivity fractions are $0.25,0.4,0.5,0.6,0.75$.
The first 20 eligible days per meter form a 1,080-task feasibility panel;
the unchanged design continues over the remaining 3,834 tasks. These are
previously used meters and forecasts, not untouched-source validation.
For each audit, only its forecasts, outer total and permitted prior-day
records enter policy construction. The shared preparation step loads
hourly endpoint tables to restore prior-day persistence; it does not
enforce physical isolation of source files. The geometry routine receives
only that audit's forecast-derived coefficients. The second adaptive query
uses those coefficients, the outer total and the first queried answer;
unqueried within-block answers enter subsequent outcome evaluation only.
This is a per-audit dependency restriction. The comparison ledgers
retain every task; none has zero initial width.

Five policies share these inputs and queries: the exact deterministic
adaptive worst-width policy; the best deterministic batch for worst width;
a randomized batch optimizing worst expected width; uniform random choice
among the ten batches; and greedy adaptive refinement. Greedy minimizes
the one-query worst width, with ties resolved by the best two-query
guarantee, then minimizes the conditional second-query worst width.
This favorable tie rule avoids an artificial gain from a poor greedy tie.
Deterministic batch ties minimize uniform-vertex mean width. The mean-optimal
batch uses the same secondary criterion on its optimal face, eliminating
pointwise-dominated optima up to numerical tolerances. It is not labeled precision-optimal.
Expectations for randomized policies sum their ten batch probabilities
exactly, rather than sampling random seeds.

\paragraph{Guarantees before retrieval.}
Two adaptive readings guarantee the primary precision in 4,661 tasks
(94.85\%). In 290 tasks (5.90\%), the guarantee is accompanied by an
adverse-profile certificate limiting \emph{every} randomized batch to
worst-case success at most 50\%. Thus these are not merely wins over a
chosen heuristic. The complete decision partition in
Table~\ref{tab:counter-guarantees} identifies when feedback is needed.
The six model pairs contribute respectively 45 (persistence/ridge),
43 (persistence/boosting), 47 (persistence/MLP), 35 (ridge/boosting),
48 (ridge/MLP) and 72 (boosting/MLP) such cases. The feasibility and
continuation panels contribute 68 and 222. The instances are identified
from permitted information before evaluating that audit's replay answers.

\begin{table}[htbp]\centering\small
\caption{Complete pre-retrieval decision partition at the unchanged primary
threshold $0.5W_0$, with two optional registers. Batch suffices means
$D\leq\tau$; feedback essential means $A\leq\tau<D$; no adaptive guarantee
means $\tau<A$ for deterministic two-query policies. All feedback-essential
cases have exact optimal randomized-batch success $1/2$. The partition is
a post-analysis diagnostic; no cases or outcomes are excluded.}
\label{tab:counter-guarantees}
\begin{tabular}{lrrrr}\toprule
Meter & Tasks & Batch suffices & Feedback essential & No adaptive guarantee\\\midrule
1 & 1,638 & 1,460 & 97 & 81\\
2 & 1,638 & 1,430 & 103 & 105\\
3 & 1,638 & 1,481 & 90 & 67\\
All & 4,914 & 4,371 & 290 & 253\\\bottomrule
\end{tabular}\end{table}

Corollary~\ref{cor:counter-feedback-triage} makes this an exhaustive
feasibility decision. A fixed batch already guarantees precision in
4,371 tasks (88.95\%). The 290 strong-separation cases exhaust the tasks
where feedback changes feasibility of a deterministic two-query guarantee;
they are not merely a favorable subset of a larger unidentified group.
In the remaining 253 tasks (5.15\%), that guarantee requires a changed
budget, threshold or information set. The supplement's recommendation
ledger returns a guaranteeing batch, an adaptive first query and conditional
second-query rule, or an unresolved flag, respectively. It uses only the
coefficient geometry and initial total. This post-analysis partition
does not replace the prespecified prevalence screen or reduce the
two-reading budget; it identifies when a sequential response is necessary
for the requested assurance.

A post-analysis executable variant uses a guaranteeing fixed batch on
those 4,371 tasks and the original adaptive policy elsewhere. This preserves
the original policy's primary precision-success indicator for every
feasible completion: where a substitution occurs, both policies guarantee
success; elsewhere they are identical. It therefore avoids the need for a
sequential response on 88.95\% of tasks without changing this precision
outcome or the two-reading count. In recorded replay, both reach 98.64\%
success; the variant's mean width is slightly higher (18.72\% versus
18.59\% of initial width). This result concerns the number of information
rounds needed by the specified protocol, not measured latency, retrieval
cost or preservation of every other objective.

\paragraph{The distribution and exact optimality of the batch bound.}
Among these 290 cases, the minimum, 25th percentile, median, 75th
percentile and maximum of the rational adverse-profile bounds are all
exactly $1/2$. All 290 equal $1/2$; none is below $1/2$, $0.4$, $0.3$
or $0.2$. This is exact equality, not rounding of nearby LP solutions.
Corollary~\ref{prop:counter-half-floor} explains why the lower tail
cannot extend below $1/2$ when a deterministic two-register policy
guarantees the precision. For each case, a constructed equal mixture of
two batches succeeds with probability at least $1/2$ on \emph{every}
completion. The independently checked finite adverse profiles give the
opposite inequality for every randomized batch. Consequently,
\[
 \underbrace{S_B^\star\geq\tfrac12}_{\text{two-batch cover of every completion}}
 \quad\text{and}\quad
 \underbrace{S_B^\star\leq\tfrac12}_{\text{finite adverse-profile certificate}}
 \quad\Longrightarrow\quad S_B^\star=\tfrac12.
\]
The finite profile set is sufficient for the upper bound; the lower bound
comes from the continuum covering proof, not from testing those profiles.
Thus the 290 cases attain the largest possible precision-success gap
between guaranteed deterministic adaptation and a randomized batch under
this two-register structure. This sharp guarantee does not assert a
similarly large gap on the recorded responses.

\begin{table}[htbp]\centering\small
\caption{Three inspectable witnesses: the earliest qualifying case from
each meter, all on 1 October 2017. P=persistence, H=boosting, R=ridge,
M=MLP; times are UTC. Query indices are hours after block start and
$t=(C_k-C_0)/E$ is the first answer. The two-branch rule is a sufficient
precision-guaranteeing policy; the measured adaptive algorithm can choose
other second queries to minimize conditional width. Each listed batch
pair is mixed with probabilities $(1/2,1/2)$. Success columns are worst-case
guarantees at $\tau=W_0/2$, not observed frequencies.}
\label{tab:counter-witnesses}
\begin{tabular}{lcllrr}\toprule
Meter: pair, block & $k$ & Second query & Batch support & $S_A$ & $S_B^\star$\\\midrule
1: P/H, 06--12 & 3 & $5$ if $t\leq1/4$, else $2$ & $\{3,5\},\{2,3\}$ & 1 & $1/2$\\
2: P/M, 06--12 & 4 & $5$ if $t\leq1/2$, else $2$ & $\{4,5\},\{2,4\}$ & 1 & $1/2$\\
3: R/M, 00--06 & 4 & $5$ if $t\leq3/4$, else $1$ & $\{4,5\},\{1,4\}$ & 1 & $1/2$\\\bottomrule
\end{tabular}\end{table}
Appendix~\ref{app:complete-counter-witness} gives one complete empirical
witness, including its rational coefficients and every batch-width vector.
All 290 exact covering intervals, batch supports, adverse weights and
case identifiers are supplied, with a program to regenerate them.
The cover and bound-distribution diagnostic was added after the primary
replay; it does not change the prespecified policies or thresholds.
The constructive mixture need not be unique or best on the recorded
trace among equally minimax-optimal mixtures. Its post-analysis descriptive replay is shown alongside the prespecified
comparisons below. The rule chooses the longest good prefix and earliest
good suffix, breaking ties by smallest second-query index, then mixes
their batches equally. It uses no internal answers. Precision optimality
holds on the 290 cases only; it does not imply best recorded performance
among all equally optimal mixtures.

For these same instances the general theorem still rules out an improvement
in the \emph{optimal worst expected width} through adaptation. The practical
consequence of the classification is therefore an objective-dependent
acquisition choice: the mean-width guarantee can use a batch, whereas
all-completions attainment of this precision requires feedback on the
certified subset. This statement concerns the uncertainty set and query
interface, not the frequency of adverse energy profiles in deployment.

\paragraph{Replay on recorded answers.}
Table~\ref{tab:counter-replay} reports the complete grid and the subset
selected by the pre-answer guarantee. On that subset, adaptive precision
success is 100\%, compared with 97.24\% for the best deterministic
worst-width batch, 80.01\% for the mean-optimal randomized batch and 96.21\%
for greedy adaptation. Across all tasks, the corresponding adaptive gains
are 0.407, 5.960 and 0.672 percentage points. All-task precision success
exceeds each nonadaptive comparator on each meter. The subset's
deterministic batch also reaches 100\% on meter 2, and its mean realized
width there is slightly smaller than adaptive; superiority at every
record or every meter/metric is not asserted.

\begin{table}[htbp]\centering\small\setlength{\tabcolsep}{4pt}
\caption{Recorded-register replay with two readings per standalone audit.
All entries are percentages: P is precision success, W is mean $W/W_0$,
and O is ordering certification (interval excluding zero). The subset is
selected by pre-answer guarantees. Randomized entries are exact expectations
over their batch mixtures. The added precision-optimal cover is defined and
proved optimal only on this subset; dashes make no full-panel claim.}
\label{tab:counter-replay}
\begin{tabular}{lrrrrrr}\toprule
 & \multicolumn{3}{c}{All 4,914 tasks} & \multicolumn{3}{c}{290-case subset}\\
Policy & P & W & O & P & W & O\\\midrule
Exact adaptive & 98.64 & 18.59 & 50.49 & 100.00 & 33.72 & 28.97\\
Best deterministic batch & 98.23 & 18.87 & 49.84 & 97.24 & 35.28 & 25.52\\
Mean-optimal random batch & 92.68 & 22.83 & 45.42 & 80.01 & 41.21 & 24.43\\
Precision-optimal cover & -- & -- & -- & 75.34 & 43.12 & 22.76\\
Uniform random batch & 63.04 & 42.59 & 24.79 & 41.79 & 53.10 & 18.76\\
Greedy adaptive & 97.96 & 18.84 & 50.33 & 96.21 & 34.67 & 28.62\\
\bottomrule\end{tabular}\end{table}

The precision-optimal cover succeeds on every recorded trace with
probability at least $1/2$: 147 cases have success one and 143 have
success $1/2$, giving $437/580=75.34\%$ overall. Its mean width is
43.12\% of initial width and its ordering-certification rate is 22.76\%.
Thus optimal worst-case precision does not entail superior recorded
success to a policy optimizing a different criterion. The matching
100\%-versus-50\% assurance result remains independent of these frequencies.

\begin{table}[htbp]\centering\small
\caption{Paired outcomes: exact adaptation versus the best deterministic
worst-width batch. A net gain is the difference between the two one-policy
columns; the rows retain countervailing cases.}\label{tab:paired-outcomes}
\begin{tabular}{llrrrr}\toprule
Tasks & Outcome & Both & Adaptive only & Batch only & Neither\\\midrule
All 4,914 & Precision & 4,817 & 30 & 10 & 57\\
All 4,914 & Ordering & 2,411 & 70 & 38 & 2,395\\
Subset 290 & Precision & 282 & 8 & 0 & 0\\
Subset 290 & Ordering & 70 & 14 & 4 & 202\\
\bottomrule\end{tabular}\end{table}

Precision success for adaptive/fixed batch is 98.84/98.72\%,
98.11/97.86\%, and 98.96/98.11\% on meters 1--3. Ordering certification
is 36.75/36.08\%, 43.65/43.77\%, and 71.06/69.66\%, respectively;
the fixed batch slightly leads on the second meter for this outcome.
On the initial 1,080 tasks, precision is 97.96/97.87\%; on the 3,834-task
continuation it is 98.83/98.33\%. These are within-source temporal summaries,
not independent replications. Repeated pairs and blocks are not treated
as independent population observations.

Normalized by squared training mean hourly energy, mean absolute MSE
widths are 0.34049 (adaptive), 0.35056 (deterministic batch), 0.42274
(mean-optimal batch), 0.80134 (uniform) and 0.34526 (greedy).
The paired ordering counts show why a precision guarantee should not be
read as an automatic model-selection decision.

\begin{table}[htbp]\centering\small
\caption{The complete predeclared precision curve. Guarantee and separation
are task counts before retrieval; strong separation uses batch upper bound
$\leq0.7$. Success is each policy's percentage on all recorded tasks.
Only $0.5W_0$ is primary.}
\label{tab:counter-curve}
\begin{tabular}{rrrrr}\toprule
$\tau/W_0$ & Guarantees & Strong separation & Adaptive success & Det.\ batch success\\\midrule
0.25 & 2,614 & 493 & 74.50 & 73.79\\
0.40 & 4,219 & 576 & 94.85 & 94.28\\
0.50 & 4,661 & 290 & 98.64 & 98.23\\
0.60 & 4,847 & 251 & 99.84 & 99.69\\
0.75 & 4,907 & 107 & 99.98 & 100.00\\\bottomrule
\end{tabular}\end{table}
\FloatBarrier

The initial screening target required strong separation in at least 10\%
of primary tasks, across all meters and at least three pairs, together
with an observed gain over every nonadaptive comparator overall and per
meter. The prevalence condition was not met: 5.90\% is reported instead,
and the primary threshold is unchanged. The other conditions hold.
The contribution is a certified, identifiable subset with a large
worst-case precision gap and a smaller observed gain, not ubiquitous
large average improvement. The complete sensitivity curve is retained
without promoting a secondary threshold to primary status.

\paragraph{Executable checks.}
The new triage recomputation independently reconstructs all ten batch
width vectors and five adaptive envelopes using rational arithmetic on
the serialized coefficients. All 4,914 finite-profile games have matching
rational primal and dual values; the largest discrepancy from saved
numerical geometry is $2.72\times10^{-16}$. Neither an adaptive nor a
fixed-batch guarantee classification is within $10^{-8}W_0$ of the primary
threshold. Thus the reported partition is unchanged by the documented
floating-point margin rule. Its ledger and executable recommendations are
checked independently of the replay outcomes.

Source audits verify the exact 91-day calendars, all 819 blocks, integer
counter telescoping and all 6,552 frozen hourly forecast rows. All 4,914
policy inputs match the original frozen forecast archive numerically,
including regenerated persistence. Independently written checks validate
each original serialized coefficient vector's adaptive envelope and
adverse-profile distribution using rational arithmetic and enumerate all
ten batches. The final second-query rule uses exact rational coefficients
and integer-counter answer ratios. This numerical hardening followed the
preliminary feasibility replay; both panels were then rerun with the same
objectives, thresholds, query dictionary and case list. The protocol and
implementation amendment are retained. The LP-derived secondary mean policy's
maximum optimal-face discrepancy is $2.15\times10^{-10}$.
Separate raw-register checks reproduce all 24,570 policy evaluations and
49,140 sharp batch intervals, with zero containment failures and zero
precision-success discrepancies. In 480 independently assembled endpoint
LPs, the largest endpoint difference is $2.23\times10^{-16}$.
Independent numerical cross-checks use scale-aware tolerances of $2\times10^{-8}$;
the exact rational witness checks apply to the serialized input instances,
not to unmodeled meter error. Online Resource 1 supplies source hashes,
policies, aggregate ledgers and the raw-source regeneration command for
the counter-dependent per-case geometry and query traces.

